\subsection{Heterogeneous fleet}
%
\rev{The fleet is generated around three nominal \aktext{vehicle} families, labeled Compact, Sport, and SUV, whose parameters are listed in Table~\ref{tab:vehicle_parameters}. These labels represent broad parameter differences rather than handling classes and are used only for fleet generation and analysis.}

\begin{table}[t]
\centering
\caption{Nominal vehicle parameters. Client-to-client parameter distributions are given in the text.}
\label{tab:vehicle_parameters}
\scriptsize
\setlength{\tabcolsep}{2.3pt}
\begin{tabular}{@{}lccc@{}}
\toprule
Parameter & Compact & Sport & SUV\\
\midrule
$m$ [kg] & 1220 & 1520 & 2550\\
$I_z$ [$\mathrm{kg\,m^2}$] & 1500 & 1800 & 5200\\
$l_f,\ l_r$ [m] & 1.05/1.45 & 1.15/1.35 & 1.30/1.55\\
$(B,C,E)_f$ & (9.0,1.27,0.93) & (13.5,1.38,0.86) & (7.5,1.22,0.95)\\
$(B,C,E)_r$ & (10.0,1.27,0.93) & (14.0,1.38,0.86) & (8.0,1.22,0.95)\\
\bottomrule
\end{tabular}
\end{table}

\begin{table}[t]
\centering
\caption{Fixed controller, actuator, and measurement-noise settings.}
\label{tab:fixed_settings}
\scriptsize
\setlength{\tabcolsep}{4pt}
\begin{tabular}{@{}ll|ll@{}}
\toprule
Quantity & Value & Quantity & Value\\
\midrule
\aktext{$K_y,K_\psi,K_r$} & $0.005\,\mathrm{m}^{-1},\ 0.1,\ 0.1\,\mathrm{s}$                         & $\sigma_X,\sigma_Y$ & $0.02\,\mathrm{m}$\\
Steering limit $\delta_{\max}$ & $0.209\,\mathrm{rad}\ (12^\circ)$   & $\sigma_\psi$ & $0.001\,\mathrm{rad}$\\
Rate limit $\dot\delta_{\max}$ & $5.236\,\mathrm{rad/s}$              & $\sigma_{v_x}$ & $0.05\,\mathrm{m/s}$\\
Road friction $\mu$ & $0.9$                                 & $\sigma_r$ & $0.001\,\mathrm{rad/s}$\\
\bottomrule
\end{tabular}
\end{table}

\rev{Each client is sampled around one of the three nominal vehicle-family anchors using three sequential heterogeneity mechanisms.
First, \aktext{vehicle-family-relative variation} is introduced by applying correlated log-normal perturbations to $m$ and $I_z$ with marginal log-standard deviation $\sigma_{\log}=0.03$ and log-space correlation $\rho=0.72$.
\akrev{The wheelbase $L$ is perturbed independently with $\sigma_{\log}=0.03$ and then split into $l_f$ and $l_r$ using the nominal-family axle-distance ratio.
Independent zero-mean Gaussian additive perturbations are applied to \(B\), \(C\), and \(E\) in each front/rear lateral and longitudinal Pacejka block, with standard deviations \(0.05\), \(0.015\), and \(0.015\), respectively.}
Second, operating-configuration variation is introduced through a payload fraction $\gamma_{\rm pay}\sim\mathcal U(0,0.25)$ placed at $x_{\rm pay}/L\sim\mathcal U(-0.20,0.25)$.
The resulting center-of-gravity shift modifies the mass, yaw inertia, and individual axle distances $l_f$ and $l_r$, while preserving the wheelbase $L=l_f+l_r$.
\akrev{Independent front and rear tire-condition factors are sampled from $\mathcal U(0.85,1.10)$ and scale the corresponding Pacejka $B$ factors in both the lateral and longitudinal tire models.}
This allows the front-to-rear balance and therefore the understeer gradient to vary within each nominal family.
\akrev{Third, independent log-normal client-to-client perturbations with $\sigma_{\log}=0.02$ are applied to $m$, $I_z$, $l_f$, $l_r$, and the front and rear Pacejka $B$ factors in both tire-force directions.
Because $l_f$ and $l_r$ are perturbed independently at this final level, the resulting wheelbase generally differs from the \aktext{vehicle-relative value}.}
These perturbations \aktext{are applied in addition to} the \aktext{vehicle-relative and operating-configuration variations} described above.
No process noise is added. The training fleet contains 150 vehicles, 50 per family.}
